% Chapter 19: Rosenbaum-Style p-Values for Matched Observational Studies
% A First Course in Causal Inference - Peng Ding
% Rewritten in Stein Style

\chapter{Rosenbaum 风格 p 值：匹配观测研究中的未测量混杂}\label{ch:19}

\section{引言：从理想实验到现实挑战}\label{sec:19-intro}

在前面的章节中，我们讨论了匹配（matching）在观测研究中的应用。匹配的核心思想是：通过构造相似的处理组和对照组单位，来近似随机化实验的条件。正如我们在第 11 章看到的，当可忽略性假设成立时，匹配可以有效地消除选择偏差。

然而，现实的观测研究往往面临一个根本性的困难：我们永远无法保证所有混淆变量都被观测到并纳入匹配。遗漏变量偏倚（omitted variable bias）是一个永恒的威胁——那些我们没有观测到的协变量，可能同时影响处理分配和结果，从而破坏因果推断的有效性。

Paul Rosenbaum 在 1987 年的开创性工作中提出了一个优雅的解决方案：不是试图证明未测量混杂不存在，而是系统地量化未测量混杂需要多大才能动摇我们的结论。\textbf{灵敏度分析}（sensitivity analysis）的核心思想是：当我们允许未测量混杂存在时，因果结论的稳健性如何？

本章专注于匹配观测研究中的 Rosenbaum 风格 p 值计算方法。我们将看到，即使在未知混杂存在的情况下，Rosenbaum 的灵敏度分析框架也能帮我们回答：这个结论在多大的未测量混杂强度下仍然是可信的？

\section{匹配观测研究的灵敏度分析模型}\label{sec:19-model}

\subsection{配对结构的记号}\label{sec:19-notation}

考虑一个匹配对照研究，我们构造了 $n$ 个配对。每个配对 $i$ 包含两个单位：$j=1$ 表示处理组，$j=0$ 表示对照组。令 $(i,1)$ 为第 $i$ 对的处理单位，$(i,0)$ 为对照单位。

每个单位 $(i,j)$ 都有两个潜在结果：$Y_{ij}(1)$ 和 $Y_{ij}(0)$。在 sharp null hypothesis 下：
\[
H_{0F}: Y_{ij}(1) = Y_{ij}(0) \quad \text{for all } i,j.
\label{eq:19-sharp-null}
\]

在理想情况下（即没有未测量混杂），每个配对中的处理和对照单位是同质的，处理分配与潜在结果独立。此时，如果我们进行随机化实验，每对中处理组接受治疗的概率应为 $\pi_{i1} = 1/2$。

\subsection{Rosenbaum 灵敏度分析模型}\label{sec:19-assumption}

然而，在观测研究中，由于未测量混杂的存在，处理分配概率可能在配对之间不相等。Rosenbaum (1987b) 提出了一个关键的灵敏度分析模型。

\begin{Assumption}[Rosenbaum 灵敏度分析模型]\label{as:19-rosenbaum}
对于第 $i$ 对中的两个单位，处理分配的 odds 比被一个有界常数 $\Gamma \geq 1$ 控制：
\[
\frac{\pi_{i1}/\pi_{i0}}{(1-\pi_{i1})/(1-\pi_{i0})} \leq \Gamma, \quad i=1,\ldots,n,
\label{eq:19-odds-ratio-bound}
\]
其中 $\pi_{i1} = \Pr\{Z_{i1}=1 \mid \text{pair } i\}$ 是配对 $i$ 中处理组接受处理的概率。
\end{Assumption}

\textbf{如何理解 $\Gamma$？} 这个参数衡量了未测量混杂导致的处理分配不均衡程度：

\begin{itemize}
  \item 当 $\Gamma = 1$ 时，\cref{eq:19-odds-ratio-bound} 意味着 $\pi_{i1} = \pi_{i0} = 1/2$，即每对中处理和对照的概率相等。这正是标准匹配实验（MPE）的情形。
  \item 当 $\Gamma > 1$ 时，处理分配概率可能在配对之间变化，变化幅度由 $\Gamma$ 控制。
  \item $\Gamma$ 越大，未测量混杂的影响越强，我们的因果结论就越脆弱。
\end{itemize}

\textbf{为什么这个模型是"自然的"？} \cref{as:19-rosenbaum} 只假设处理分配的 odds 比有界，而没有对潜在的未测量混杂做任何具体假设。这种半参数化的设定使得灵敏度分析具有普遍的适用性——无论未测量混杂的具体形式如何，只要其导致的 odds 比不超过 $\Gamma$，我们的 p 值计算就是有效的。

\section{最坏情况 p 值}\label{sec:19-worst-case}

\subsection{检验统计量}\label{sec:19-test-stat}

在 sharp null hypothesis $H_{0F}$ 下，我们观察到的是处理效应为零的"无效"数据。利用这个性质，我们可以构造 Fisher randomization test 风格的检验统计量。

令 $q_i$ 为第 $i$ 对的某种处理效应度量，常见的 choices 包括：

\begin{enumerate}
  \item \textbf{Sign statistic}：$q_i = Z_{i1} - Z_{i0}$（即配对中处理 vs 对照的指示函数）
  \item \textbf{Pair $t$ statistic}：$q_i = Y_{i1} - Y_{i0}$（观测到的处理组与对照组结果差）
  \item \textbf{Wilcoxon signed-rank statistic}：基于 $|Y_{i1} - Y_{i0}|$ 的符号秩统计量
\end{enumerate}

检验统计量定义为：
\[
T = \sum_{i=1}^n S_i q_i,
\label{eq:19-test-stat}
\]
其中 $S_i \in \{0,1\}$ 是配对 $i$ 的处理指示（$S_i=1$ 表示该对参与检验），$q_i$ 是配对 $i$ 的贡献权重。

在标准的完全随机化检验中，$S_i$ 在零假设下是已知的（或可精确计算）。但在灵敏度分析框架下，$S_i$ 的分布受到 \cref{as:19-rosenbaum} 的约束。

\subsection{最坏情况分布}\label{sec:19-worst-case-dist}

\begin{Theorem}[最坏情况 p 值]\label{th:19-worst-case}
在 \cref{as:19-rosenbaum} 下，零假设 $H_{0F}$ 成立时，配对 $i$ 的处理指示 $S_i$ 服从最坏情况的 Bernoulli 分布：
\[
S_i \sim \text{Bernoulli}\left(\frac{\Gamma}{1+\Gamma}\right),
\label{eq:19-s_i-dist}
\]
且各配对的 $S_i$ 之间相互独立。
\end{Theorem}

\textbf{为什么是最坏情况？} 给定 odds ratio 上界 $\Gamma$，$S_i$ 的概率可以在 $[\pi_{\min}, \pi_{\max}]$ 之间变化，其中
\[
\pi_{\max} = \frac{\Gamma}{1+\Gamma}, \quad \pi_{\min} = \frac{1}{1+\Gamma}.
\]
为了得到最大的 p 值（即最不利于拒绝零假设的情况），我们选择 $S_i$ 的期望为 $\pi_{\max}$。这对应于处理效应"看起来最小"的那种分布。

\cref{th:19-worst-case} 的证明表明，在所有满足 \cref{as:19-rosenbaum} 的分布中，\cref{eq:19-s_i-dist} 给出了最保守的 p 值。\footnote{定理 \cref{th:19-worst-case} 的完整证明见附录 \cref{sec:appendix-19-worst-case-proof}。}

\subsection{p 值的计算}\label{sec:19-p-value}

基于 \cref{th:19-worst-case}，我们可以计算最坏情况 p 值。对于给定的观测检验统计量 $T^{\mathrm{obs}}$，双边 p 值为：
\[
\text{p}_{\Gamma} = \Pr_{\Gamma}\left(|T - T^{\mathrm{obs}}| \geq |T^{\mathrm{obs}}|\right),
\label{eq:19-worst-p-value}
\]
其中 $T$ 服从 \cref{eq:19-s_i-dist} 下 $S_i$ 的混合分布。

当 $q_i \geq 0$ 时，我们可以简化计算。定义
\[
T = \sum_{i=1}^n S_i q_i, \quad S_i \sim \text{Bernoulli}\left(\frac{\Gamma}{1+\Gamma}\right).
\]
由于 $q_i$ 是固定的（由观测数据决定），$T$ 的分布是 $n$ 个独立 Bernoulli 随机变量的加权求和。

\section{实例分析}\label{sec:19-examples}

\subsection{LaLonde 数据重访}\label{sec:19-lalonde}

LaLonde (1986) 的数据集是因果推断领域最著名的实例之一。该数据集包含一个工作培训项目（National Supported Work Demonstration）的参与者及其结果。

我们考虑使用检验统计量：
\[
T = \sum_{i=1}^n S_i |\hat{\tau}_i|,
\label{eq:19-lalonde-stat}
\]
其中 $\hat{\tau}_i$ 是第 $i$ 对的处理效应估计。

\textbf{灵敏度分析结果}：

\begin{itemize}
  \item 在理想 MPE（$\Gamma = 1$）下，模拟得到的 p 值为 $0.002$，强烈拒绝零假设。
  \item 当 $\Gamma = 1.1$ 时（仅允许 $10\%$ 的处理分配 odds 偏离），最坏情况 p 值上升到 $0.011$。
  \item 当 $\Gamma = 1.3$ 时，最坏情况 p 值超过 $0.05$，在常规显著性水平下无法拒绝零假设。
\end{itemize}

更精确地说，$\Gamma = 1.233$ 是临界值：当未测量混杂导致的 odds 比不超过 $1.233$ 时，我们仍能在 $0.05$ 水平下拒绝零假设；超过这个值，结论就不再稳健。

\begin{Example}[LaLonde 数据的灵敏度曲线]\label{ex:19-lalonde-sensitivity}
图 \ref{fig:19-lalonde-sensitivity} 展示了 p 值随 $\Gamma$ 变化的灵敏度曲线。曲线的陡峭程度告诉我们结论对未测量混杂的敏感程度——曲线越平缓，结论越稳健。
\end{Example}

\cref{fig:19-lalonde-sensitivity} 由以下 R 代码生成。首先使用 \texttt{Matching} 包构造匹配数据集，然后使用 \texttt{sensitivitymw} 包中的 \texttt{senmw} 函数计算不同 $\Gamma$ 值下的最坏情况 p 值：

\begin{lstlisting}[language=R]
library("Matching")
library("sensitivitymvm")
library("sensitivitymw")
dat <- read.table("cps1re74.csv", header = TRUE)
dat$u74 <- as.numeric(dat$re74 == 0)
dat$u75 <- as.numeric(dat$re75 == 0)
y = dat$re78
z = dat$treat
x = as.matrix(dat[, c("age", "educ", "black",
                      "hispan", "married", "nodegree",
                      "re74", "re75", "u74", "u75")])

## 1-1 配对匹配
matchest = Match(Y = y, Tr = z, X = x)
ytreated = y[matchest$index.treated]
ycontrol = y[matchest$index.control]
datamatched = cbind(ytreated, ycontrol)

## 计算 pair t 统计量 tau_i
taumatched = (ytreated - ycontrol)

## 不同 Gamma 下的 p 值序列
Gamma = seq(1, 1.4, 0.001)
Pvalue = sapply(Gamma, function(g) senmw(taumatched, gamma = g, method = "t")$pval)

## 找到 Gamma_star（使 p-value = 0.05 的临界值）
gammastar = Gamma[which(Pvalue >= 0.05)[1]]
gammastar  # 1.233
\end{lstlisting}

图 \ref{fig:19-lalonde-sensitivity} 的左图展示了 $\hat{\tau}_i$ 的直方图，右图展示了 p 值随 $\Gamma$ 变化的曲线，其中虚线标记了 $\Gamma = 1.233$ 的临界位置。

\begin{figure}[ht]
  \centering
  \safeincludegraphics[width=0.9\textwidth]{R/figures/ch19-lalonde-sensitivity.pdf}
  \caption{LaLonde 数据的灵敏度分析：左图为配对处理效应 $\hat{\tau}_i$ 的分布；右图为 p 值随 $\Gamma$ 变化的曲线，虚线 $\Gamma = 1.233$ 标记了 $0.05$ 显著性水平的临界值。}
  \label{fig:19-lalonde-sensitivity}
\end{figure}

\subsection{Rosenbaum 包中的其他实例}\label{sec:19-rosenbaum-packages}

Rosenbaum 的 R 包（如 \texttt{sensitivity2x2} 和 \texttt{sensitivitymv}）提供了便捷的灵敏度分析工具。我们介绍两个应用场景。

\textbf{实例 1：焊接工人 DNA 洗脱率研究（erpcp 数据）}。该数据集来自 \texttt{sensitivitymw} 包，包含 $n=39$ 个配对（焊工与对照），基于年龄和吸烟协变量匹配。结局是 DNA 洗脱率。

\begin{lstlisting}[language=R]
library("sensitivitymw")
data(erpcp)
taue = erpcp[, 1] - erpcp[, 2]  # 配对内差异

## 图 19.3(a): 左图直方图 + 右图灵敏度曲线
Gamma = seq(1, 5, 0.005)
Pvalue = sapply(Gamma, function(g) senmw(taue, gamma = g, method = "t")$pval)
gammastar = Gamma[which(Pvalue >= 0.05)[1]]
## gammastar = 3.366
\end{lstlisting}

\textbf{实例 2：吸烟者血铅水平研究（lead250 数据）}。该数据集来自 \texttt{sensitivitymw} 包，包含 $n=250$ 个配对（每日吸烟者与非吸烟者对照），基于性别、年龄、种族、教育程度和家庭收入进行匹配。结局是血铅水平（$\mu$g/L）。

\begin{lstlisting}[language=R]
data(lead250)
taul = lead250[, 1] - lead250[, 2]  # 配对内差异

Gamma = seq(1, 2.5, 0.001)
Pvalue = sapply(Gamma, function(g) senmw(taul, gamma = g, method = "t")$pval)
gammastar = Gamma[which(Pvalue >= 0.05)[1]]
## gammastar = 1.476
\end{lstlisting}

图 \ref{fig:19-rosenbaum-packages} 展示了这两个数据集的分析结果。erpcp 数据要求 $\Gamma$ 超过 $3.37$ 才能使 p 值超过 $0.05$，而 lead250 数据仅要求 $\Gamma > 1.48$，说明后者对未测量混杂更敏感。

\begin{figure}[ht]
  \centering
  \safeincludegraphics[width=0.9\textwidth]{R/figures/ch19-rosenbaum-packages.pdf}
  \caption{(a) erpcp 数据：焊工 vs 对照的 DNA 洗脱率灵敏度分析；(b) lead250 数据：吸烟者 vs 非吸烟者的血铅水平灵敏度分析。左列为 $\hat{\tau}_i$ 直方图；右列为 p 值随 $\Gamma$ 变化的曲线。}
  \label{fig:19-rosenbaum-packages}
\end{figure}

\section{习题}\label{sec:19-homework}

\begin{HomeworkProblem}[A model for Assumption \ref{assume::rosenbaum-sensitivity}]\label{exr:19-1}

Assumption \ref{assume::rosenbaum-sensitivity} has the following equivalent form.


\begin{assumption}
\label{assume::rosenbaum-equivalent}
The propensity score satisfies the following model
\[
\log \frac{  \pi_{ij} }{  1-\pi_{ij}  } = g(X_i) + \gamma U_{ij} ,\quad (i=1,\ldots, n; j=1,2)
\]
where $g(\cdot)$ is an unknown function and $U_{ij} \in [0, 1]$ is a a bounded unobserved covariate for unit $(i, j)$.
\end{assumption}


Show that if Assumption \ref{assume::rosenbaum-equivalent} holds, then Assumption \ref{assume::rosenbaum-sensitivity} must hold with $\Gamma = e^\gamma$.


Remark: \citet[][page 108]{rosenbaum2002design} also shows that if Assumption \ref{assume::rosenbaum-sensitivity} holds, then Assumption \ref{assume::rosenbaum-equivalent} must hold for some $U_{ij}$'s. Assumption \ref{assume::rosenbaum-equivalent} may be easier to interpret because $\gamma$ measures the log of the conditional odds ratio of $U$ on the treatment; see Chapter \ref{sec::logistic-regression} for the interpretation of the coefficients in the logistic regression.



\end{HomeworkProblem}

\begin{HomeworkProblem}[Application of Rosenbaum's approach]\label{exr:19-2}
Re-analyze Example \ref{eg::chan-bmi-ATE} using Rosenbaum's approach based on matching.



\end{HomeworkProblem}

\begin{HomeworkProblem}[Recommended reading]\label{exr:19-3}


\citet{rosenbaum2015two} provides a tutorial for his two \ri{R} packages for sensitivity analysis with matched observational studies.


\end{HomeworkProblem}
